\chapter{Frontend React}\label{cap:frontend}

\section{Estructura}

\begin{lstlisting}[caption={Layout del frontend}]
fakenews-frontend/src/
|-- App.jsx, main.jsx
|-- components/        # Componentes reutilizables (UI, dashboard,
|                      # admin, auth, landing, common)
|-- hooks/             # Hooks reutilizables
|-- lib/               # Utilidades genericas, i18n, constantes
|-- locales/{en,es}/   # JSONs de traduccion
|-- pages/             # Vistas de ruta
|-- services/          # Capa de servicios (auth, billing, analysis,
|                      # supabase, history, account, admin)
|-- store/             # Stores Zustand
`-- test/              # Configuracion Vitest
\end{lstlisting}

\section{Enrutamiento y vistas}

React Router v6 organiza las rutas en públicas (\texttt{/},
\texttt{/login}, \texttt{/register}, \texttt{/forgot-password},
\texttt{/reset-password}) y autenticadas (\texttt{/dashboard/*},
\texttt{/admin/*}). Una capa de protección consulta
\texttt{authStore} para redirigir a \texttt{/login} a usuarios sin
sesión, y a \texttt{/dashboard} a usuarios sin permisos
administrativos que intenten acceder al panel.

\section{Internacionalización}

i18next con detección automática de idioma del navegador y persistencia
en \texttt{localStorage}. Cambio en caliente desde el menú de
preferencias. Traducciones organizadas en \texttt{namespaces}
(\texttt{auth.json}, \texttt{dashboard.json}, \texttt{landing.json}).

\section{Sistema de diseño}

Tailwind CSS configurado con tokens propios de marca (colores
primarios, tipografía). Componentes shadcn/ui generados con CLI
(\path{components.json}) sobre Radix UI.

\section{Análisis y verificación}

El \emph{hook} \path{src/hooks/useAnalysisFlow.js} centraliza el flujo
de análisis: validación del input, llamada al servicio
\path{src/services/analysis.js} y actualización del store. La función
\texttt{verifyClaims} llama a \texttt{POST /verify} y, ante 403,
muestra un mensaje localizado indicando que la funcionalidad sólo está
disponible en el plan \emph{Super Pro}.

El componente
\path{src/components/dashboard/VerificationReport.jsx} renderiza el
veredicto global, la lista de \emph{claims} con su confianza y los
enlaces a evidencias (\texttt{rel="noopener noreferrer"}).

\section{Pruebas}

Vitest + React Testing Library cubren componentes críticos
(\path{src/pages/__tests__/}, \path{src/lib/__tests__/}). Los tests
mockean los servicios para aislar la capa UI/estado.
